\documentclass[UTF8,twocolumn]{ctexart}

\usepackage[a4paper,margin=1.8cm]{geometry}
\usepackage{amsmath,amssymb,bm}
\usepackage{booktabs}
\usepackage{multirow}
\usepackage{graphicx}
\usepackage{array}
\usepackage{caption}
\usepackage{hyperref}
\hypersetup{colorlinks=true,linkcolor=blue,urlcolor=blue,citecolor=blue}

\setlength{\columnsep}{0.7cm}
\setlength{\parindent}{1.2em}

\newcommand{\placeholderfig}[2]{%
\begin{figure}[t]
\centering
\fbox{\parbox[c][#2][c]{0.92\columnwidth}{\centering 【图片占位符：#1】\\（此处留空，便于后续插入原文图像）}}
\caption{#1}
\end{figure}
}

\newcommand{\placeholderfigstar}[2]{%
\begin{figure*}[t]
\centering
\fbox{\parbox[c][#2][c]{0.95\textwidth}{\centering 【图片占位符：#1】\\（此处留空，便于后续插入原文图像）}}
\caption{#1}
\end{figure*}
}

\title{\bfseries XFall：具有交叉模态监督（Cross-Modal Supervision）的域自适应（Domain Adaptive）Wi\!-\!Fi（Wi-Fi）跌倒检测（Fall Detection）}
\author{Guoxuan Chi（IEEE Member）, Guidong Zhang（IEEE Member）, Xuan Ding（IEEE Member）,\\
Qiang Ma（IEEE Member）, Zheng Yang（IEEE Fellow，通讯作者）,\\
Zhenguo Du, Houfei Xiao, Zhuang Liu}
\date{}

\begin{document}
\maketitle

\begin{abstract}
近年来，人类跌倒检测（fall detection）系统的需求持续增长。在现有方法中，基于 Wi\!-\!Fi（Wi-Fi）的跌倒检测因其普适性（pervasiveness）而成为最有前景的方案之一。然而，当应用到新域（domain）时，现有基于 Wi\!-\!Fi 的方案由于泛化能力（generalizability）不足，会出现严重的性能下降。本文提出 XFall——一种基于 Wi\!-\!Fi 的域自适应（domain-adaptive）跌倒检测系统。XFall 从三个方面解决泛化问题：为推进跨环境（cross-environment）感知，XFall 利用一种与室内布局和设备部署无关的环境无关特征——速度分布谱（speed distribution profile, SDP）；为保证对所有跌倒类型（fall types）的灵敏性，设计了一种基于注意力（attention）的编码器，通过关联输入的空间与时间维度来提取通用的跌倒表示；为在有限 Wi\!-\!Fi 数据下训练大模型，设计了交叉模态（cross-modal）学习框架，在训练过程中采用预训练视觉模型（pre-trained visual model）进行监督。我们在最新商用无线产品之一上实现并评估了 XFall，并在真实场景中进行了为期一年的部署。结果表明：XFall 总体准确率（accuracy）为 96.8\%，漏报率（miss alarm rate, MAR）为 3.1\%，误报率（false alarm rate, FAR）为 3.3\%，在域内（in-domain）与跨域（cross-domain）评估中均优于现有最先进（state-of-the-art）方案。
\end{abstract}

\noindent\textbf{索引词：}域适配（domain adaptation），跌倒检测（fall detection），统计电场（statistical electric field），变压器编码器（transformer encoder），交叉模态监督（cross-modal supervision）。

\section{引言（Introduction）}
人类跌倒（human falls）是重大的公共健康问题。据估计，每年约有 684,000 起致死性跌倒（fatal falls），是全球非故意伤害致死（unintentional injury death）的第二大原因\cite{ref1}。过去十年，大量研究致力于降低意外跌倒带来的死亡与伤害。

尽管跌倒检测方案不断涌现，包括视觉（vision-based）系统、可穿戴设备（wearable devices）与传感器（sensor-based）方案，但各自存在局限。视觉系统\cite{ref2,ref3,ref4}受隐私（privacy）与相机视场（field-of-view, FoV）限制，且在非视距（non-line-of-sight, NLoS）条件下无效。可穿戴设备\cite{ref5,ref6}虽具有创新性，但因侵入性（intrusive）往往面临用户接受度（user acceptance）问题。依赖先进雷达（radar）\cite{ref7}或声学传感器（acoustic sensors）\cite{ref8}的方案则成本较高，阻碍了广泛部署。因此，亟需一种普及、可靠且非侵入式的室内跌倒检测方法，以弥补现有方案的关键缺口。

过去数十年，Wi\!-\!Fi 基础设施（infrastructures）已广泛部署。作为充斥室内空间的普遍无线信号，Wi\!-\!Fi 已被用于构建多种非接触式（contactless）感知系统\cite{ref9,ref10,ref11,ref12,ref13}。早期研究\cite{ref14,ref15,ref16}表明，Wi\!-\!Fi 信道状态信息（Channel State Information, CSI）可用于跌倒检测。但现有方法的关键挑战在于：当环境超出其训练域时，跨域泛化（cross-domain generalization）能力不足，从而限制了实际部署。要实现“训练一次、到处使用（train once, use anywhere）”的一体化 Wi\!-\!Fi 跌倒检测系统，我们面临三大挑战：

\textbf{挑战 1：如何提取环境无关（environment-independent）的特征。}
以往系统要么使用原始 CSI 幅度/相位（raw CSI amplitude and phase）\cite{ref14,ref17}等原始特征，要么使用由 CSI 生成的时频图（spectrogram）特征，例如多普勒频谱图（Doppler Frequency Spectrogram, DFS）\cite{ref15}与离散小波变换谱（Discrete Wavelet Transformation, DWT）\cite{ref18}。但由于多径效应（multipath effect），原始 CSI 模式在不同室内结构间差异显著，从而掩盖由人体活动引起的特征；DFS 与 DWT 虽能反映与室内结构相对独立的环境动态（environment dynamics），却高度依赖设备部署与用户位置\cite{ref19}。因此，上述特征都难以直接适配不同环境。

\textbf{挑战 2：如何设计模型学习通用的人体跌倒表示（general human fall representation）。}
跌倒是复杂活动，不同跌倒类型会在时空上产生不同信号扰动。现有方案使用经典机器学习模型，如支持向量机（support vector machine, SVM）\cite{ref15}，或深度神经网络（deep neural networks），如多层感知机（multilayer perceptron, MLP）与卷积神经网络（convolutional neural network, CNN）\cite{ref7}。但这些方法普遍未显式建模输入特征的空间维与时间维关联。例如，有些人跌倒时会扶住家具或靠墙，导致信号波动并不显著，这类跌倒常难以被现有模型检测。

\textbf{挑战 3：如何在 Wi\!-\!Fi 数据有限时训练大模型。}
要获得良好泛化性，通常需要拥有数千万参数的深度模型，但这要求大量标注（labeled）的 Wi\!-\!Fi 感知数据进行训练。遗憾的是，Wi\!-\!Fi 数据采集与标注极为费力，尤其需要在线标注（on-line labeling）\footnote{在线标注（on-line labeling）指在采集数据的同时进行标注。}，因为 Wi\!-\!Fi 感知数据难以通过人类观察直接解释\cite{ref20}。因此，收集足量标注数据是巨大挑战。

为应对上述挑战，我们设计并实现了 XFall——一种基于 Wi\!-\!Fi 的域自适应跌倒检测系统。为得到环境无关无线特征，我们深入电磁场传播（electromagnetic field propagation）的基本原理：受统计电磁场理论（statistical electromagnetic field theory）启发，构建速度分布谱（SDP），可适配不同室内结构、设备部署与用户位置。为挖掘通用跌倒表示，我们提出时空注意力变压器编码器（Spatio-Temporal-Attention-based Transformer Encoder, STATE），学习输入特征在空间与时间维度的关联与依赖，从而在更高抽象层次刻画连续活动。为在有限标注数据下训练大模型，我们提出交叉模态统一特征学习（Cross-modal Unified Representation Learning, CURL）框架，借鉴知识蒸馏（knowledge distillation）思想\cite{ref21}：利用视觉域（visual domain）的跌倒检测网络在训练阶段对 STATE 进行监督，从而降低对标注 Wi\!-\!Fi 数据的需求。

我们在商用 Wi\!-\!Fi 产品华为 AX3 Pro 上实现 XFall，并在 70 个真实设置中进行了为期一年的评估，涵盖多样室内布局、接入点（access point, AP）部署与用户位置。我们与三种最先进（SOTA）Wi\!-\!Fi 跌倒检测系统对比：DeFall\cite{ref16}、TLFall\cite{ref18} 与 FallDeFi\cite{ref15}。结果显示，XFall 平均 MAR 为 3.1\%，平均 FAR 为 3.3\%，总体准确率为 96.8\%，相较 DeFall、TLFall、FallDeFi 分别提升 14.5\%、22.3\%、23.1\%，体现了 XFall 的真实场景优势。

本文贡献如下：
\begin{itemize}
\item 提出 XFall：首个基于 Wi\!-\!Fi 的域自适应跌倒检测系统。该设计可适配不同环境（室内布局、AP 部署与用户位置不同）并对所有跌倒类型保持灵敏，是迈向实用与泛在（ubiquitous）Wi\!-\!Fi 感知的重要一步。
\item XFall 中提出的环境无关特征与时空编码器具有独特优势，可直接用于其他 Wi\!-\!Fi 感知任务（如手势识别 gesture recognition、步态识别 gait recognition）以提升泛化能力。
\item 提出的 CURL 框架有效降低标注 Wi\!-\!Fi 数据需求，使 Wi\!-\!Fi 感知系统具备小样本学习（few-shot learning）能力。
\item 在 AX3 Pro 上实现并评估 XFall，这是基于支持 802.11ax（Wi\!-\!Fi 6）标准的最新商用无线产品构建真实无线感知应用的先行尝试。
\end{itemize}

本文其余部分组织如下：第 II 节回顾相关工作；第 III 节给出系统概览；第 IV 节详述 SDP；第 V 节详述 STATE；第 VI 节介绍 CURL；第 VII 节给出实现与评估；第 VIII 节总结。

\section{相关工作（Related Work）}
本节简要总结最相关研究。

\subsection{基于 CSI 的跌倒检测（CSI-based fall detection）}
信道状态信息（CSI）在 Wi\!-\!Fi 感知中被广泛利用，尤其在跌倒检测中展示了商用设备捕获细粒度运动模式的潜力。WiFall\cite{ref14} 基于 CSI 幅度实现跌倒检测；在此基础上，RTFall\cite{ref17} 同时利用 CSI 幅度与不同天线间相位差检测跌倒事件。FallDeFi\cite{ref15} 对 CSI 幅度做时频分析（time-frequency analysis），并提出顺序前向选择（sequential forward selection）算法筛选对环境变化更鲁棒的特征。TLFall\cite{ref18} 提取 DWT 轮廓并采用迁移学习（transfer learning）降低环境差异影响。DeFall\cite{ref16} 是开创性工作，从原始 CSI 中提取环境无关速度特征，且无需密集训练。FallViewer\cite{ref22} 提出一系列 CSI 校准（calibration）算法减少环境成分、增强跌倒活动影响。

上述方法通常从特征中提取统计特性，并用经典机器学习（如 SVM、随机森林 random forest）区分跌倒与非跌倒。尽管贡献显著，但“无需重新配置即可跨环境无缝适配”的目标仍未达成。与上述方法相比，XFall 引入 SDP——一种对环境与部署差异不敏感（agnostic to domain variations）的特征，使系统在变化条件下更稳健；同时借助 STATE 识别跌倒相关的高层特征（high-level features），提升检测准确率。

\subsection{其他模态的跌倒检测（Fall detection based on other modalities）}
除 CSI 外，跌倒检测还探索了视觉、可穿戴设备与专用雷达等技术。视觉方法通过 RGB 相机、事件相机（event camera）与深度相机（depth camera）采集图像并提取人体运动\cite{ref2,ref3,ref4,ref23}。可穿戴方法通常依赖惯性测量单元（Inertial Measurement Unit, IMU），包含加速度计（accelerometer）、陀螺仪（gyroscope）与倾角计（inclinometer）等\cite{ref5,ref6,ref24,ref25,ref26}。专用雷达方面，RF-Fall\cite{ref7} 通过调频连续波雷达（Frequency Modulated Continuous Wave, FMCW）提取空间热力图并用 CNN 识别跌倒；Acoustic-FADE\cite{ref8} 使用环形麦克风阵列识别跌倒。受统计无线感知启发，VeCare\cite{ref27} 首次提出统计声学感知（statistical acoustic sensing），可用音频信号做速度估计。本研究通过利用普遍存在的 Wi\!-\!Fi 信号构建非侵入式检测系统，填补了“泛在且易融入日常环境”的技术缺口。

\subsection{基于 CSI 的活动识别（CSI-based activity recognition）}
CSI 的应用不仅限于跌倒检测，还涵盖手势、步态识别与被动跟踪等。E-eyes\cite{ref9} 率先用商用 Wi\!-\!Fi 区分原地活动。WiFiU\cite{ref28}、Indotrack\cite{ref29} 与 Widar\cite{ref30} 分别基于 DFS 做步态识别与被动跟踪。WiSpeed\cite{ref31} 首次引入统计电磁（statistical EM）方法，通过 CSI 功率自相关函数（auto-correlation function, ACF）估计速度，并在 GaitWay\cite{ref32} 中扩展到 CSI 的 ACF。Widar3.0\cite{ref19} 提取体坐标速度轮廓（body-coordinate velocity profile, BVP）。EI\cite{ref33} 与 CrossSense\cite{ref34} 分别结合对抗网络（adversarial networks）与迁移学习实现活动感知。SLNet\cite{ref35} 提出面向射频信号（RF signals）的复值神经网络（complex-valued neural network），通过频谱分析与深度学习协同设计，展示了对多种任务的通用性。近期 RF-Diffusion\cite{ref13} 通过数据增强提升感知系统准确性与鲁棒性。但将这些洞见用于跌倒检测，仍需克服真实场景中环境与情境（situational）变化带来的挑战。XFall 通过提取域无关特征，实现无需域特定训练（domain-specific training）的精确检测，拓展了跌倒检测技术边界。

\section{系统概览（System Overview）}
XFall 是基于现成商用（off-the-shelf）Wi\!-\!Fi 设备的人体跌倒检测系统。图 1 展示了系统架构：在监测区域周围部署一对 Wi\!-\!Fi 设备以采集 CSI 流（CSI streams）；同时采集 RGB 相机视频流，用于交叉模态训练（cross-modal training）。XFall 包含三部分：速度分布谱（SDP）生成、时空注意力变压器编码器（STATE）与交叉模态统一表示学习（CURL）框架。系统接收跌倒与正常活动（normal activities）的 CSI 序列后执行 SDP 生成算法：先建立统计电场模型（statistical electric field model），再计算 CSI 自相关函数以生成 SDP，包含环境无关的人体运动信息。随后将 SDP 序列输入 STATE：它在每一时刻提取高层空间特征，并在整个序列上提取时间特征；注意力模块促使网络聚焦于与跌倒对应的时空特征，从而提升分类精度。STATE 输出可提取通用跌倒表示（general fall representation, GFR）。训练阶段，CURL 引入预训练视觉模型监督 STATE，实现跨模态学习。

需要强调的是：CURL 仅在训练阶段工作；测试阶段 XFall 仅需 Wi\!-\!Fi 数据，不需要任何视觉输入。

\placeholderfigstar{图 1：XFall 系统架构。}{4.2cm}

\section{速度分布谱（Speed Distribution Profile, SDP）}
已有研究\cite{ref36}表明，人全身各部位的独特速度分布（velocity distribution）是活动识别（activity recognition）最有效指标之一。在无线特征（如衰减 attenuation、到达角 Angle of Arrival, AoA、到达时间 Time of Flight, ToF）中，DFS 轮廓包含环境速度信息，被广泛用于活动识别\cite{ref37}。但 DFS 是环境相关（environment-dependent）的：与收发端部署与用户位置/朝向高度相关，难以跨域应用\cite{ref19}。本节提出域无关特征——速度分布谱（SDP），可适配不同室内布局、用户位置与设备部署。我们先介绍统计电场模型基础理论（第 IV-A 节），在此基础上给出从 CSI 表达电场（第 IV-B 节）以及 SDP 的计算方法（第 IV-C 节）。

\subsection{统计电场模型（Statistical Electric Field Model）}
如图 2 左侧所示，大多数室内空间可建模为富散射（rich scattering）环境：散射体（scatterers）包括人体部位、家具与墙体等，可视为漫反射（diffusive），能将无线信号向各方向反射\cite{ref31}。发射机（transmitter, Tx）发出的连续电磁波（electromagnetic, EM waves）经多次散射后由接收机（receiver, Rx）接收。

通常，室内 EM 波可由其电场（electric field）完全表征。令 $\vec{E}_{Rx}(t,f)$ 表示在时间 $t$、频率 $f$ 处接收的电场。为分析环境空间动态，将其分解为
\begin{equation}
\vec{E}_{Rx}(t,f)=\sum_{\alpha\in\Omega_s(t)}\vec{E}_{\alpha}(t,f)+\sum_{\beta\in\Omega_d(t)}\vec{E}_{\beta}(t,f),
\label{eq1}
\end{equation}
其中 $\Omega_s(t)$ 与 $\Omega_d(t)$ 分别为静态与动态散射体集合。

进一步分析单个分量 $\vec{E}_{\alpha}(t,f)$（由第 $\alpha$ 个散射体散射的 EM 波）。可将其解释为来自各方向的入射（incident / incoming）EM 波的积分（见图 2 右侧）。对仰角（elevation）$\phi$ 与方位角（azimuth）$\theta$ 的入射波，定义其角谱（angular spectrum）为 $\vec{A}(\phi,\theta)$，波数向量（wave-number vector）为
$\vec{k}=-\frac{2\pi f}{c}(\sin\phi\cos\theta,\sin\phi\sin\theta,\cos\phi)^{T}$，
其中 $c$ 为光速。设散射体 $\alpha$ 的速度为 $\vec{v}_{\alpha}$，由麦克斯韦方程（Maxwell's equations）\cite{ref38}可得
\begin{equation}
\vec{E}_{\alpha}(t,f)=\int_{0}^{2\pi}\int_{0}^{\pi}\vec{A}(\phi,\theta)\exp\!\left(-j\,\vec{k}\cdot\vec{v}_{\alpha}t\right)\sin\phi\,d\phi\,d\theta,
\label{eq2}
\end{equation}
其中 $z$ 轴与散射体 $\alpha$ 的运动方向对齐。式 \eqref{eq2} 建立了速度 $v_{\alpha}$ 与接收电场之间的联系。

然而，缺乏室内结构先验时，传播过程难以解析：既无法确定散射体位置，也难以确定 $\vec{A}(\phi,\theta)$。因此我们不构建确定性（deterministic）模型，而是将室内空间视为混响腔（reverberation cavity）\cite{ref38}，建立电磁场的统计模型（statistical model）。具体而言，$\vec{E}_{\alpha}(t,f)$ 被视为来自均匀分布方向、具有随机相位与极化（polarizations）的大量独立波的叠加。因此可得到两点结论：（1）角谱 $\vec{A}(\phi,\theta)$ 服从圆对称高斯随机变量（circularly symmetric Gaussian random variable）\cite{ref39}；（2）对任意 $\alpha,\beta\in\Omega_d$ 与任意 $t_1,t_2$，入射波分量 $\vec{E}_{\alpha}(t_1,f)$ 与 $\vec{E}_{\beta}(t_2,f)$ 不相关（uncorrelated）。

据此，$\vec{E}_{\alpha}(t,f)$ 可近似为平稳随机过程（stationary random process），其自相关函数（ACF）为
\begin{equation}
\rho_{\vec{E}_{\alpha}}(\tau,f)=
\frac{\langle \vec{E}_{\alpha}(t,f),\vec{E}_{\alpha}(t+\tau,f)\rangle}{
\sqrt{\|\vec{E}_{\alpha}(t,f)\|^{2}\,\|\vec{E}_{\alpha}(t+\tau,f)\|^{2}}},
\label{eq3}
\end{equation}
其中 $\langle\cdot\rangle$ 为内积算子，$\tau$ 为时延（time lag）。已有研究给出详细证明\cite{ref40}：
\begin{equation}
\langle \vec{E}_{\alpha}(t,f),\vec{E}_{\alpha}(t+\tau,f)\rangle
=E_{\alpha}^{2}(f)\,\frac{\sin(k v_{\alpha}\tau)}{k v_{\alpha}\tau},
\label{eq4}
\end{equation}
其中 $E_{\alpha}^{2}(f)$ 定义为 $\vec{E}_{\alpha}(t,f)$ 的功率（power）。聚合所有独立动态散射分量，可得接收电场的 ACF：
\begin{equation}
\rho_{\vec{E}_{Rx}}(\tau,f)=
\frac{1}{\sum_{\beta\in\Omega_d}E_{\beta}^{2}(f)}
\sum_{\alpha\in\Omega_d}E_{\alpha}^{2}(f)\,\frac{\sin(k v_{\alpha}\tau)}{k v_{\alpha}\tau}.
\label{eq5}
\end{equation}
式 \eqref{eq5} 建立了接收电场与环境动态（即各散射体速度 $v_{\alpha}$）之间的联系：通过分析接收电场的 ACF，可推断环境的速度分布（speed distribution）。

\placeholderfig{图 2：富散射环境示意：发射波在到达接收端前被环境扩散散射。}{3.7cm}

\subsection{从 CSI 到散射体速度（From CSI to Scatter Speed）}
在商用 Wi\!-\!Fi 设备上获取接收电场是关键挑战。令 $X(t,f)$ 与 $Y(t,f)$ 分别为频率为 $f$ 的子载波（subcarrier）在时间 $t$ 的发射与接收信号波。注意到 $Y(t,f)=\|\vec{E}_{Rx}\|$。根据 CSI 定义，$H(t,f)=Y(t,f)/X(t,f)$。与式 \eqref{eq1} 类似，可将 CSI 分解为
\begin{equation}
H(t,f)=\sum_{\alpha\in\Omega_s(t)}H_{\alpha}(t,f)+\sum_{\beta\in\Omega_d(t)}H_{\beta}(t,f),
\label{eq6}
\end{equation}
其中 $H_{\alpha}(t,f)$ 为对应散射体 $\alpha$ 的 CSI 分量。由于 $X(t,f)$ 是 802.11 协议中预定义的长训练字段（long training field），可视为常数，因此 CSI 幅度 $|H(t,f)|$ 与电场范数成正比：
\begin{equation}
|H(t,f)|\propto \|\vec{E}_{Rx}(t,f)\|,\qquad
|H_{\alpha}(t,f)|\propto \|\vec{E}_{\alpha}(t,f)\|.
\label{eq7}
\end{equation}
根据式 \eqref{eq4} 与式 \eqref{eq5}，电场协方差（covariance）与自相关均退化为与电场方向无关的标量。因此可将式 \eqref{eq7} 代入式 \eqref{eq4} 与式 \eqref{eq5}，建立 CSI 与环境速度分布之间的统计关系：
\begin{equation}
\rho_{H}(\tau,f)=
\frac{\mathrm{cov}(H(t,f),H(t+\tau,f))}{\mathrm{cov}(H(t,f),H(t,f))}
=\sum_{\alpha\in\Omega_d} w_{\alpha}\,\mathrm{sinc}(k v_{\alpha}\tau),
\label{eq8}
\end{equation}
其中 $\mathrm{cov}(\cdot,\cdot)$ 表示协方差，$\sigma_{\alpha}(t,f)$ 表示 CSI 分量幅度 $|H_{\alpha}(t,f)|$ 的标准差（standard deviation），$w_{\alpha}$ 为权重。式 \eqref{eq8} 表明：CSI 的 ACF 可视作散射体速度 $v_{\alpha}$ 经非线性滤波 $\mathrm{sinc}(\cdot)$ 后的加权和。

\subsection{生成速度分布谱（Generate the Speed Distribution Profile）}
实际中，Wi\!-\!Fi 网卡（NIC）上报的 CSI 数据形成二维复矩阵 $H\in\mathbb{C}^{N_T\times N_S}$，其中 $N_T$ 为包（packets）数、$N_S$ 为子载波数。$H(t_i,f_j)$ 表示第 $i$ 个包、第 $j$ 个子载波的 CSI 复样本。

一个 SDP $S$ 由三项关键参数决定：CSI 样本数 $N_T$、时延分辨率（time lag resolution）$\Delta T$ 与时延样本数 $N_{\Delta}$。每个 SDP 由三维 ACF 张量（tensor）$\rho\in\mathbb{R}^{N_{\Delta}\times W_T\times N_S}$ 生成，其中元素
\begin{equation}
\rho(n,i,j)=\rho_H(n\Delta T,f_j)=
\frac{|H(t_i,f_j)H^{*}(t_i-n\Delta T,f_j)|}{
|H(t_i,f_j)|\,|H(t_i-n\Delta T,f_j)|},
\label{eq9}
\end{equation}
$H^{*}$ 为共轭（conjugate）。

在此基础上，对所有子载波合并，并对 ACF 张量每一列做概率归一化（probabilistic normalization），得到 $S\in\mathbb{R}^{N_{\Delta}\times W_T}$：
\begin{equation}
S(n,i)=
\frac{\sum_{j=1}^{N_S}\rho(n,i,j)/N_S}{
\sum_{n=1}^{N_{\Delta}}\rho(n,i,j)}.
\label{eq10}
\end{equation}
为便于理解，基于式 \eqref{eq8} 定义随机过程
$P_u(\tau,v)=\sum_{\alpha\in\Omega_d}w_{\alpha}\mathrm{sinc}(k v_{\alpha}\tau)$。如图 3 所示，$S$ 的每一列表示在不同 $\tau$ 处对 $u$ 的离散采样，包含环境的速度分布。因此，可将 SDP 视为矩阵，其两维分别包含环境动态（environmental dynamics）的空间与时间信息。

我们进一步分析 SDP 提取算法的复杂度：式 \eqref{eq9} 与 \eqref{eq10} 的逐元素（element-wise）操作作用于三维张量，单次操作的时间复杂度为 $O(N_{\Delta}\times W_T\times N_S)$。实际中三维规模相近，因此总复杂度约为 $2\times O(N_{\Delta}\times W_T\times N_S)\approx O(N^3)$，具有现实可用性。

\placeholderfig{图 3：速度分布谱（SDP）示意：二维分别对应速度的空间与时间分布。}{3.0cm}

\subsection{与其他特征的比较（Comparison with Different Features）}
为更直观展示 SDP 的优势，我们将其与两种特征——相位差（phase difference）与 DFS——在不同域中对比。实验在客厅与会议室中进行（见图 8）：志愿者在不同朝向与位置下执行几乎相同的跌倒动作，Wi\!-\!Fi 设备记录 CSI。相位差按 RTFall\cite{ref17} 计算；DFS 通过短时傅里叶变换（short-time Fourier transform, STFT）获得，并按 FallDeFi\cite{ref15} 采用带通滤波（band-pass filter）以降噪；SDP 按第 IV-C 节方法提取。

图 4 展示了跌倒过程中三类特征的模式。可见，CSI 相位差对朝向、位置与环境高度敏感，随域变化剧烈波动，难以表征活动。DFS 在相同位置但不同朝向（location \#1、orientation \#1/\#2）、或相同朝向但不同位置（orientation \#1、location \#1/\#2）时呈现不同频移；在不同环境下也出现不同模式。理论上，DFS 仅能捕获目标径向速度（radial velocity），难以感知切向运动（tangential motion），因此依赖设备部署与用户位置/朝向。相比之下，SDP 在不同朝向、位置与环境下对同一跌倒类型呈现近似一致特征，体现出强域泛化能力（domain generalization）。

\placeholderfigstar{图 4：不同域设置（环境、位置、朝向）下跌倒的相位差、DFS 与 SDP。}{5.8cm}

\section{空间与时间注意力网络（Spatial and Temporal Attention Network）}
在 XFall 中，我们提出图 5 所示的时空注意力变压器编码器（STATE），用于从输入 SDP 中充分挖掘通用跌倒表示（GFR）。STATE 的基本构件是变压器编码器（transformer encoder）\cite{ref41}（第 V-A 节），用于学习输入特征的关联与依赖。受其结构启发，我们分别设计空间变压器编码器（spatial transformer encoder，第 V-B 节）与时间变压器编码器（temporal transformer encoder，第 V-C 节），分别沿 SDP 的两维提取特征。

如图 5 所示，首先将 SDP 切分为 $W_T$ 个列向量 $S_i$（大小 $N_{\Delta}\times 1$），表示在时刻 $t_i$ 的空间速度分布；序列 $\{S_i\}$ 输入空间编码器，提取环境速度的空间信息。随后将所有空间嵌入（spatial embeddings）汇聚形成序列数据，输入时间编码器建模时序特征。

\placeholderfigstar{图 5：时空注意力变压器编码器（STATE）结构示意。}{5.6cm}

\subsection{变压器编码器（Transformer Encoder）}
图 6 展示了一个变压器编码器块的通用结构。变压器编码器是基于自注意力（self-attention）的模型，能有效从视频流（video streams）与文本序列（text sequences）等输入嵌入中提取高层描述\cite{ref41}。其过程包括：对输入序列进行线性投影（linear projection）与位置嵌入（position embedding）；多头自注意力（multi-head self-attention）捕获输入块（patches）间依赖；再通过全连接前馈网络（feed-forward blocks）进一步建模。每个块后加入残差连接（residual connection）与层归一化（layer normalization），提升网络敏感性并避免梯度消失（vanishing gradient）问题\cite{ref42}。

多头自注意力的目标是建立不同时间块之间的交互。给定输入 $X\in\mathbb{R}^{L\times d_{in}}$，线性投影生成查询（queries）$Q\in\mathbb{R}^{L\times d_Q}$、键（keys）$K\in\mathbb{R}^{L\times d_K}$ 与值（values）$V\in\mathbb{R}^{L\times d_V}$：
\begin{equation}
Q=XW_Q,\quad K=XW_K,\quad V=XW_V,
\label{eq11}
\end{equation}
其中 $W_Q\in\mathbb{R}^{d_{in}\times d_Q}$、$W_K\in\mathbb{R}^{d_{in}\times d_K}$、$W_V\in\mathbb{R}^{d_{in}\times d_V}$ 为投影参数，且 $d_Q=d_K$。注意力输出为值的加权和，权重由查询与键的点积（dot-product）确定：
\begin{equation}
\mathrm{Attention}(Q,K,V)=\mathrm{softmax}\!\left(\frac{QK^{T}}{\sqrt{d_K}}\right)V.
\label{eq12}
\end{equation}
实践中使用 $h$ 个注意力头（attention heads）的多头注意力，独立计算 $h$ 次以捕获不同子空间信息。将各头输出拼接并再投影得到最终输出：
\begin{equation}
\mathrm{MultiHead}(Q,K,V)=(a_1,a_2,\ldots,a_h)W_O,
\label{eq13}
\end{equation}
其中 $a_i=\mathrm{Attention}(XW_Q^{i},XW_K^{i},XW_V^{i})$，$W_O\in\mathbb{R}^{h d_V\times d_O}$ 为最终投影参数。为保证输入输出维度一致，设 $d_{in}=h d_V$。通过多头自注意力块，可获得输入序列的深层交互信息（deep interactions）。

\placeholderfig{图 6：变压器编码器结构示意。}{3.2cm}

\subsection{空间特征提取（Spatial Feature Extraction）}
直接将基础变压器编码器用于无线信号的空间特征提取会遇到两点问题：其一，编码器生成的输出向量数量大，计算开销高，且用于分类任务可能导致训练集过拟合（overfitting）；其二，尽管变压器不强制指定块位置，但无线信号中输入块的顺序通常对应不同空间状态下的功率分布（power distribution），具有重要物理含义。

为解决上述问题，我们在基础编码器上加入分类标记（classification token, \texttt{CLS}）与可学习位置嵌入（learnable position embeddings），构建空间变压器编码器。空间编码器关注空间速度分布而非时间变化。其输入为列向量 $S_i$（$N_{\Delta}\times 1$），描述与空间速度分布相关的随机过程。为加速训练与推理，将 $S_i$ 重塑为二维矩阵 $X_i\in\mathbb{R}^{L\times d_S}$，满足 $N_{\Delta}=L\times d_S$。重塑后每一行可视为一个输入块，共有 $L$ 个块。为提升分类学习效果，将 $\texttt{CLS}\in\mathbb{R}^{1\times d_S}$ 插入到块序列开头，作为整段输入的表示\cite{ref43}；变压器对 \texttt{CLS} 的输出状态用作分类特征。于是编码器最终输入为
\[
X_i'=(\texttt{CLS};X_i)\in\mathbb{R}^{(L+1)\times d_S}.
\]
同时加入位置嵌入 $P_i\in\mathbb{R}^{(L+1)\times d_S}$ 以保留位置信息，得到 $\tilde{X}'_i=X'_i+P_i$。经空间编码器得到高层特征 $Y_i\in\mathbb{R}^{(L+1)\times d_S}$，取 \texttt{CLS} 对应输出 $y_i^{0}\in\mathbb{R}^{1\times d_S}$ 作为空间编码器输出。所有 $W_T$ 个空间编码器的输出序列将作为后续时间建模的输入。

对每一时刻而言，空间编码器通过自注意力自适应地对不同功率分布分配注意力，更关注跌倒相关的速度空间分布，而较少关注其他速度成分。

\subsection{时间特征提取（Temporal Feature Extraction）}
本节介绍时间编码器，用于提取时间特征。与空间编码器类似，时间编码器也会遇到输出向量多与输入块位置缺失的问题。为此我们提出时间注意力编码器，使其更关注速度分布剧烈变化的位置（moment）。

时间注意力网络的输入为 $W_T$ 个空间编码器输出，记为
\[
y^{0}=(y^{0}_{1},y^{0}_{2},\ldots,y^{0}_{W_T})\in\mathbb{R}^{W_T\times d_S},
\]
其中每个 $y^{0}_{i}$ 可视为一个输入块，共 $W_T$ 个块。与第 V-B 节相同，我们加入 \texttt{CLS} 与位置嵌入，将基础变压器用于时间特征提取。经过时间变压器编码器后，提取 \texttt{CLS} 的输出 $z^{0}\in\mathbb{R}^{1\times d_S}$ 作为时间编码器输出。随后通过若干全连接层（fully-connected layers）将 $z^{0}$ 映射为 GFR（general fall representation），表示人体跌倒的通用表示。

为展示时间注意力机制的效果，志愿者在同一设置下执行五种跌倒类型（slip、trip、lose-balance、kneel-then-fall、walk-then-fall）。Wi\!-\!Fi 设备采集 CSI 序列，按第 IV-C 节提取 SDP，并由时间注意力编码器计算各时刻的注意力权重。图 7 展示五类跌倒在不同时刻的 SDP 与注意力权重。可见不同跌倒类型的 SDP 差异较大；时间注意力为每个时刻分配不同权重。在跌倒过程中，编码器会对更高速度时刻给予更高注意力，而高速度是判断跌倒的关键要素。借助注意力机制，时间编码器聚焦于跌倒过程中的关键运动，从而捕获本质时间信息。

\placeholderfigstar{图 7：不同跌倒类型（slip、trip、lose-balance、kneel-then-fall、walk-then-fall）的 SDP 与时间注意力。}{5.2cm}

\section{交叉模态统一表示学习（Cross-Modal Unified Representation Learning, CURL）}
STATE 能有效提取人体活动的时空表示（spatio-temporal representation）。但变压器模型参数量大，通常需要大量训练数据。Wi\!-\!Fi 数据集的采集与标注费力，且需要在线标注（on-line labeling），不同于图像或文本数据。

受知识蒸馏（knowledge distillation）\cite{ref21}启发，我们提出交叉模态统一表示学习（CURL）框架：利用与 Wi\!-\!Fi 数据同步采集的视觉信号，将视觉域的分类能力迁移到无线特征域（wireless feature domain），从而在训练阶段提升学习效率。CURL 对 XFall 的收益体现在两方面：首先，视觉分类网络的特征图（feature maps）用于监督 STATE 的训练，可有效降低对标注 Wi\!-\!Fi 数据的需求；其次，CURL 促使 STATE 学习视觉表示与 Wi\!-\!Fi 特征之间的关联，提升系统精度并加速训练。

需要指出：视觉网络仅用于在训练阶段“引导（guide）”STATE；一旦训练完成，XFall 仅以 Wi\!-\!Fi 信号为输入。

\subsection{视觉网络（Vision-based network）}
为实现有效交叉模态学习，我们选择一种简单且有效的模型 C3D\cite{ref44}，其由若干 3D 卷积层（3D convolution layers）与全连接层构成。构建视觉跌倒检测网络时，我们先选择预训练的 C3D 模型，并用采集的人体跌倒与正常活动视频进行微调（fine-tune）。微调过程中冻结卷积层参数，仅调整全连接层。得到训练好的 C3D 后，从全连接层中部提取特征图 $r_e$，用于后续交叉模态学习。

\subsection{交叉模态监督（Cross-modal supervision）}
训练阶段，将同步视频与 Wi\!-\!Fi 数据输入 CURL 框架，共同训练 STATE 与用于分类的 MLP。基于\cite{ref45}，同时学习网络特征图与最终分类结果可获得更好性能。因此，STATE 有两个训练目标：（1）STATE 输出的 GFR 应匹配视觉模型输出特征图 $r_e$；（2）最终跌倒检测结果应匹配真实标签（ground truth）。

因此损失函数（loss function）由两部分组成。一方面，采用均方误差（mean square error, MSE）评价 GFR 的监督效果。设第 $i$ 个样本中，STATE 输出为 $r_i$，视觉监督为 $r_{ei}$，监督损失为
\begin{equation}
L_{sv}=\sum_{i=1}^{N}\|r_{ei}-r_i\|^{2}/N,
\label{eq14}
\end{equation}
其中 $N$ 为训练样本总数。另一方面，采用二元交叉熵（binary cross-entropy）评价分类性能：
\begin{equation}
L_c=-\sum_{i=1}^{N_0}\left(Y_i\log P(Y_i')+(1-Y_i)\log(1-P(Y_i'))\right),
\label{eq15}
\end{equation}
其中 $Y_i$ 与 $Y_i'$ 分别为第 $i$ 个样本的真实标签与预测标签。最终总损失为
\begin{equation}
L=L_c+\lambda L_{sv},
\label{eq16}
\end{equation}
$\lambda$ 为平衡系数（balance factor）。通过反向传播（backpropagation）与优化算法（optimization algorithms）训练 STATE 与分类 MLP 层。

\section{评估（Evaluation）}
\subsection{实验方法（Experiment Methodology）}
\subsubsection{实现（Implementation）}
我们在商用 Wi\!-\!Fi 产品华为 AX3 Pro 上实现 XFall 原型，该产品支持最新 802.11ax 标准。发射端与接收端中心频率为 5.825\,GHz；二者均配备两根天线，形成 $2\times 2$ 多输入多输出阵列（multiple-input multiple-output, MIMO）。发射端以 350 包/秒的注入速率（injection rate）发送 Wi\!-\!Fi 数据包以提取 CSI。软件实现采用混合编程（hybrid-programming）：使用 MATLAB 进行信号处理与 SDP 生成，使用 PyTorch 构建与训练深度模型。

\subsubsection{评估设置（Evaluation setup）}
表 I 给出了评估配置与方法，旨在反映真实可用性（real-world applicability）。实验覆盖三种典型室内环境：办公室、客厅与小会议室，并形成超过 70 种不同域（domains），包括不同室内布局、AP 部署与用户位置（见图 8）。如此广泛的测试条件有助于确保评估能覆盖实际部署中可能遇到的多样场景。具体地：办公室中固定发射端位置，将接收端放置在 16 个不同位置；客厅与会议室中发射端与接收端均放置在给定位置。对每种接收端部署，设置 2--4 个志愿者可能跌倒的位置。使用 1080P RGB 相机以 20 FPS 记录地面真实（ground truth），并与 CSI 采集同步，从而精确对齐跌倒事件与数据。

我们的实验不仅验证了 XFall 在多种真实场景下的功能性，也展示了其与标准商用 Wi\!-\!Fi 设备的兼容性。XFall 仅需单条 Wi\!-\!Fi 链路（single Wi-Fi link）即可有效工作，与家庭使用场景高度契合。实验设置与部署方式也得到了具有丰富真实部署经验的无线设备厂商认可，进一步佐证了系统的可行性。

\subsubsection{数据采集（Data collection）}
五名身高与体型（somatotypes）各异的志愿者参与实验。采集跌倒样本时，志愿者需在图 8 指定位置执行不同跌倒类型的跌倒动作；采集正常样本时，志愿者在实验场景中执行日常动作。共采集 1,000+ 跌倒样本与 2,800+ 正常样本。跌倒类型包括：滑倒（slip）、绊倒（trip）、失衡（lose-balance）、跪下再跌倒（kneel-then-fall）、走动中跌倒（walk-then-fall）。正常活动包括：慢走、快走、坐下、站起、跳跃、弯腰拾取（bend-and-pickup）与深蹲（squatting）。所有实验均通过机构审查委员会（Institutional Review Board, IRB）批准。

\subsubsection{指标（Metrics）}
参考\cite{ref15}，我们使用两项简单指标评估性能：漏报率（Missed Alarm Rate, MAR）与误报率（False Alarm Rate, FAR）。MAR 衡量系统对跌倒的敏感性（sensitivity），定义为被错误漏检的跌倒事件占比；FAR 衡量系统避免误报的精确性（precision），定义为被错误识别为跌倒的非跌倒事件占比。

\begin{table*}[t]
\centering
\caption{表 I：详细评估设置（Detailed Evaluation Setup）}
\label{tab1}
\begin{tabular}{@{}p{3.0cm}p{12.8cm}@{}}
\toprule
方面（Aspect） & 细节（Details）\\
\midrule
\multirow{5}{*}{设备配置（Device Config）}
& 设备类型（Device Type）：Huawei AX3 Pro\\
& 标准（Standard）：IEEE 802.11ax\\
& 频率（Frequency）：5.825\,GHz\\
& 天线（Antenna）：$2\times 2$ MIMO\\
& 包速率（Packet Rate）：350 packets per second\\
\midrule
\multirow{7}{*}{数据采集（Data Collection）}
& 总样本（Total Samples）：1,000 fall samples，2,800 normal activity samples\\
& 环境（Environments）：Office，Living Room，Meeting Room\\
& 志愿者（Volunteers）：5 volunteers with diverse heights and body types\\
& 跌倒类型（Fall Types）：Slip，Trip，Lose-balance，Kneel-fall，Walk-fall\\
& 正常活动（Normal Activities）：Walking，Sitting，Standing，Jumping，Bending，Squatting\\
& 用户状态（User States）：each scenario 2--4 different locations and orientations\\
& AP 部署（AP Deployment）：18 different Tx-Rx settings\\
\midrule
地面真实（Groundtruth）
& 1080P RGB camera recorded at 20 FPS\\
\midrule
评估（Evaluation）
& 指标（Metrics）：MAR，FAR，F1-Score\\
& 实验（Experiments）：Cross-type，Cross-environment，Cross-user\\
\bottomrule
\end{tabular}
\end{table*}

\placeholderfigstar{图 8：三种评估环境布局。}{5.0cm}

\subsection{总体性能（Overall Performance）}
我们将 XFall 与三种最先进（SOTA）跌倒检测方案对比：FallDeFi\cite{ref15}、TLFall\cite{ref18} 与 DeFall\cite{ref16}。FallDeFi 与 TLFall 分别提取 DFS 与 DWT 轮廓并选取若干统计特征，再用 SVM 分类跌倒与正常活动；DeFall 从原始 CSI 提取速度序列（speed streams），并用模板匹配（template matching）分类。

为公平比较，我们将全部数据划分为训练集与测试集，并采用十折交叉验证（ten-fold cross-validation）。结果如图 9：XFall 总体准确率 96.8\%，平均 MAR 3.1\%，平均 FAR 3.3\%，优于现有工作。

为评估鲁棒性（robustness），我们进一步进行跨类型（cross-type）与跨环境（cross-environment）对比：测试集包含训练集中未出现的跌倒类型或环境。如图 10 与图 11 所示，XFall 在不同域上均保持更高性能，体现域自适应能力（domain-adaptive performance）。

根据业内经验基准（industry-informed empirical benchmarks），在跨类型与跨环境条件下将 MAR 与 FAR 控制在 10\% 以下，对商用跌倒检测系统的可用性至关重要。现有方案（如 FallDeFi、TLFall）往往难以满足该建议，DeFall 也仅在一致条件下勉强对齐；而 XFall 在多种域上均能满足这一广泛认可的准则，强调了其真实场景实用性（practical utility）。

与 XFall 相比，FallDeFi 与 TLFall 的分类模型（SVM）更简单，难以充分利用输入特征信息；DeFall 使用估计的单一速度值作为输入，虽能反映跌倒，但不同活动的速度可能相近，模板匹配在分类上存在局限。XFall 提取的是速度分布（speed distribution），比单一速度值更能描述活动；同时，时空注意力模型也帮助 XFall 获得更佳性能。

\placeholderfigstar{图 9：与 SOTA 的总体对比。}{3.8cm}
\placeholderfigstar{图 10：跨类型（cross-type）与 SOTA 对比。}{3.8cm}
\placeholderfigstar{图 11：跨环境（cross-environment）与 SOTA 对比。}{3.8cm}

\subsection{鲁棒性分析（Robustness Analysis）}
\subsubsection{不同跌倒类型下的性能（Performance across different fall types）}
我们采集五类典型跌倒样本：slip、trip、lose-balance、kneel-then-fall、walk-then-fall，分别标为 \#1--\#5。选取一种跌倒类型与部分正常样本作为测试集，其余四种跌倒类型与剩余正常样本作为训练集。图 12 显示：XFall 在不同跌倒类型上保持一致性能，表明其能将已学习的跌倒知识迁移到未直接训练过的跌倒类型上，显著提升泛化能力，使其更适用于真实场景。该结果的原因主要有二：（1）SDP 聚焦时空速度分布而非单一速度值，可捕获更全面的跌倒模式；（2）STATE 通过注意力机制更充分地分析跌倒活动，提取更细致的跌倒相关信息。因此，XFall 能适配不同跌倒类型。

\subsubsection{不同环境下的性能（Performance across different environments）}
评估 XFall 在不同环境下的适配性对理解实际应用范围至关重要。我们仅用图 8 中环境 \#a 数据训练 XFall，并在三个环境（含两个未见环境）中测试。图 13 显示：跨环境时性能变化很小；即使迁移到新环境，MAR 与 FAR 仍约为 7.0\%。这种稳定性体现了 XFall 的跨环境鲁棒性与强泛化潜力。其根本原因在于 SDP 的设计：SDP 聚焦于与环境无关的人体跌倒普遍特征，从而有效隔离环境特定噪声（environment-specific noise）。因此，XFall 在任意环境训练后，都能在未见环境中准确检测跌倒，体现其适配性与部署就绪性（readiness for real-world deployment）。

\subsubsection{不同用户下的性能（Performance across different users）}
为评估跨个体泛化能力，我们让五名志愿者（\#1--\#5）参与实验。训练时使用其中四人组合的数据，测试时使用剩余未见用户数据。图 14 显示：XFall 在不同用户上表现一致，MAR 与 FAR 平均约 4.0\%，证明其能有效识别与用户无关的鲁棒跌倒特征（user-independent robust fall characteristics）。

\subsection{微基准（Micro Benchmarks）}
\subsubsection{特征的影响（Impact of features）}
我们比较四类不同抽象层次的特征：去噪 CSI（denoised CSI）、DFS 轮廓、DWT 轮廓与 SDP，并将其输入同一深度模型（第 V 节）进行训练与测试。训练/测试集均由三种环境样本构成。图 15 显示：以 MAR+FAR 衡量，SDP 相较去噪 CSI、DFS、DWT 平均分别降低 33.8\%、5.8\%、8.7\%。与 CSI 相比，SDP 利用电磁场先验挖掘具有物理意义的速度分布信息，使深度模型更易学习与优化；DFS 与 DWT 同时包含环境与活动信息，而 SDP 更关注与环境/位置/朝向无关的活动特异信息，因此实践中性能更佳。结果表明 SDP 是跌倒检测的理想特征。

\subsubsection{学习模型的影响（Impact of learning models）}
我们评估 STATE 分类器的有效性。以 SDP 为输入，分别用 CNN、长短期记忆网络（Long Short-Term Memory, LSTM）与 STATE 分类跌倒与正常活动。图 16 显示：STATE 相较 CNN 与 LSTM 的 MAR+FAR 平均分别降低 4.0\% 与 5.7\%。这表明本模型能分别从空间与时间维度挖掘活动信息；注意力机制还能聚焦关键时刻与关键空间特征。因此，相比传统模型，STATE 更能利用人体动作信息，获得更好性能。

\subsubsection{视频监督的有效性（Effectiveness of video supervision）}
我们通过改变训练样本数，比较有无视频监督（video supervision）时系统性能，并用 F1 分数（F1-score）作为评估指标（综合 MAR 与 FAR）。图 17 显示：有监督系统的 F1 分数始终高于无监督系统；换言之，在达到同等 F1 分数时，有监督系统所需训练样本更少。随着训练样本增多，两者差距逐渐减小。结论是：视觉网络监督能提升深度模型训练效率并减少所需 Wi\!-\!Fi 训练数据量。

\subsubsection{采样率影响（Impact of sampling rates）}
前述实验采样率为 350\,Hz。为评估不同采样率下性能，我们将 CSI 下采样（downsample）到 100--300\,Hz，并相应调整 STATE 的输入形状。图 18 显示：采样率从 350\,Hz 降到 200\,Hz 时，MAR 与 FAR 分别仅上升 1.7\% 与 1.6\%；进一步降到 100\,Hz 时，MAR 与 FAR 分别上升 2.7\% 与 4.0\%。理论上跌倒动作较快，高采样率更敏感。结果显示采样率高于 200\,Hz 后性能提升不显著，因此在现成商用设备（commercial off-the-shelf, COTS）上无需极高采样率也可获得理想性能。

\subsubsection{窗口大小影响（Impact of window size）}
经验上跌倒持续时间约 1000--1500\,ms。我们评估不同时间窗口（1500/1250/1000/750/500\,ms）对性能影响。图 19 显示：窗口从 1500\,ms 降到 1000\,ms 时，MAR 与 FAR 分别仅上升 2.1\% 与 4.6\%；但进一步降到 500\,ms 时，MAR 与 FAR 分别上升 17.9\% 与 10.2\%。理论上窗口过短会导致跌倒活动捕获不完整，从而性能下降。结果表明应设置 1500\,ms 以获得最佳性能。

\subsubsection{系统时延分析（System latency analysis）}
为验证效率，我们在低功耗商用笔记本上测试端到端时延。图 20 显示：以 1500\,ms CSI 为输入，平均端到端时延为 59.9\,ms，其中 SDP 计算平均 46.8\,ms、STATE 分类平均 13.1\,ms。系统时延主要来自 SDP 计算，尤其是 ACF 计算。结论：XFall 可实现实时监测（real-time monitoring），适用于对即时性要求高的场景。

\subsubsection{模型可扩展性分析（Model scalability analysis）}
为理解 STATE 在不同配置下的表现，我们测试 9 种配置：注意力块数量（attention block numbers）32B/48B/64B 与注意力头数（attention heads）4H/8H/16H 的组合（见图 21）。结果表明：增加块数与头数会提高参数量与性能，但收益递减。例如参数从 10M 增到 40M，性能提升超过 5\%；从 40M 增到 160M，性能仅提升约 3\%。此外，增加头数在提升性能与参数量的同时，并未显著增加总体计算负载，这归因于头数增加会降低每条注意力路径的计算需求。总体而言，STATE 具备良好可扩展性（scalability），进一步增加块与头有望提升精度，体现 XFall 的潜力；同时可在不同计算平台上灵活部署（versatile）。

\subsubsection{多 AP 协作的有效性（Effectiveness of multiple AP collaboration）}
虽然 XFall 设计为在单条 Wi\!-\!Fi 链路上有效工作，但我们也探索多链路协作以揭示性能上限。我们在同一环境部署 1--3 套 XFall 单元，并采用两种经典集成学习（ensemble learning）策略：XGBoost\cite{ref46} 与随机森林（Random Forest）\cite{ref47} 融合输出。图 22 显示：融合 2 个 AP 后，XGBoost 与随机森林将准确率提升至 98.3\% 与 98.1\%；扩展到 3 个 AP 可提升至 98.7\% 与 98.4\%，但计算复杂度增加超过 2 倍。XGBoost 整体更优且提升更稳定。综上，多 AP 协作可显著突破 XFall 固有性能上界，是未来增强方向。

\placeholderfigstar{图 12：不同跌倒类型下的性能。}{3.6cm}
\placeholderfigstar{图 13：不同环境下的性能。}{3.6cm}
\placeholderfigstar{图 14：不同用户下的性能。}{3.6cm}
\placeholderfigstar{图 15：不同特征对性能的影响。}{3.6cm}
\placeholderfigstar{图 16：不同学习模型对性能的影响。}{3.6cm}
\placeholderfigstar{图 17：视频监督有效性。}{3.6cm}
\placeholderfigstar{图 18：采样率影响。}{3.6cm}
\placeholderfigstar{图 19：窗口大小影响。}{3.6cm}
\placeholderfigstar{图 20：系统时延分析。}{3.6cm}
\placeholderfigstar{图 21：模型可扩展性分析。}{3.8cm}
\placeholderfigstar{图 22：多 AP 协作有效性。}{3.8cm}

\section{结论（Conclusion）}
本文提出并实现 XFall——首个基于 Wi\!-\!Fi 的域自适应跌倒检测系统。XFall 构建环境无关特征速度分布谱（SDP），并采用新型时空注意力编码器（STATE）学习通用跌倒表示（GFR）。借助交叉模态统一表示学习（CURL）框架，引入视觉模型监督训练过程，从而有效降低对标注 Wi\!-\!Fi 数据的需求。评估结果显示，XFall 在域内与跨域评估中均优于现有最先进 Wi\!-\!Fi 跌倒检测方案，是迈向实用与泛在 Wi\!-\!Fi 感知的重要进展。

\section*{参考文献（References）}
\begin{thebibliography}{99}

\bibitem{ref1} (2021) Falls. [Online]. Available: \url{https://www.who.int/news-room/fact-sheets/detail/falls}
\bibitem{ref2} G. Mastorakis and D. Makris, “Fall detection system using kinect’s infrared sensor,” \emph{Journal of Real-Time Image Processing}, 2014.
\bibitem{ref3} C. Krupitzer, T. Sztyler, J. Edinger, M. Breitbach, H. Stuckenschmidt, and C. Becker, “Hips do lie! a position-aware mobile fall detection system,” in \emph{2018 IEEE International Conference on Pervasive Computing and Communications (PerCom)}, IEEE, 2018.
\bibitem{ref4} Z.-P. Bian, J. Hou, L.-P. Chau, and N. Magnenat-Thalmann, “Fall detection based on body part tracking using a depth camera,” \emph{IEEE journal of biomedical and health informatics}, 2014.
\bibitem{ref5} L.-J. Kau and C.-S. Chen, “A smart phone-based pocket fall accident detection, positioning, and rescue system,” \emph{IEEE journal of biomedical and health informatics}, 2014.
\bibitem{ref6} T. Theodoridis, V. Solachidis, N. Vretos, and P. Daras, “Human fall detection from acceleration measurements using a recurrent neural network,” in \emph{International conference on biomedical and health informatics}. Springer, 2017, pp. 145–149.
\bibitem{ref7} Y. Tian, G.-H. Lee, H. He, C.-Y. Hsu, and D. Katabi, “Rf-based fall monitoring using convolutional neural networks,” \emph{Proceedings of the ACM on Interactive, Mobile, Wearable and Ubiquitous Technologies}, 2018.
\bibitem{ref8} Y. Li, K. Ho, and M. Popescu, “A microphone array system for automatic fall detection,” \emph{IEEE Transactions on Biomedical Engineering}, 2012.
\bibitem{ref9} Y. Wang, J. Liu, Y. Chen, M. Gruteser, J. Yang, and H. Liu, “E-eyes: device-free location-oriented activity identification using fine-grained wifi signatures,” in \emph{Proceedings of the 20th annual international conference on Mobile computing and networking}, 2014, pp. 617–628.
\bibitem{ref10} W. Wang, A. X. Liu, M. Shahzad, K. Ling, and S. Lu, “Understanding and modeling of wifi signal based human activity recognition,” in \emph{Proceedings of the 21st annual international conference on mobile computing and networking}, 2015, pp. 65–76.
\bibitem{ref11} G. Chi, Z. Yang, J. Xu, C. Wu, J. Zhang, J. Liang, and Y. Liu, “Wi-drone: wi-fi-based 6-dof tracking for indoor drone flight control,” in \emph{Proceedings of the 20th Annual International Conference on Mobile Systems, Applications and Services}, 2022, pp. 56–68.
\bibitem{ref12} Y. Gao, G. Chi, G. Zhang, and Z. Yang, “Wi-prox: Proximity estimation of non-directly connected devices via sim2real transfer learning,” in \emph{GLOBECOM 2023-2023 IEEE Global Communications Conference}. IEEE, 2023, pp. 5629–5634.
\bibitem{ref13} G. Chi, Z. Yang, C. Wu, J. Xu, Y. Gao, Y. Liu, and T. X. Han, “Rf-diffusion: Radio signal generation via time-frequency diffusion,” \emph{arXiv preprint arXiv:2404.09140}, 2024.
\bibitem{ref14} Y. Wang, K. Wu, and L. M. Ni, “Wifall: Device-free fall detection by wireless networks,” \emph{IEEE Transactions on Mobile Computing}, vol. 16, no. 2, pp. 581–594, 2016.
\bibitem{ref15} S. Palipana, D. Rojas, P. Agrawal, and D. Pesch, “Falldefi: Ubiquitous fall detection using commodity wi-fi devices,” \emph{Proceedings of the ACM on Interactive, Mobile, Wearable and Ubiquitous Technologies}, vol. 1, no. 4, pp. 1–25, 2018.
\bibitem{ref16} Y. Hu, F. Zhang, C. Wu, B. Wang, and K. R. Liu, “Defall: Environment-independent passive fall detection using wifi,” \emph{IEEE Internet of Things Journal}, vol. 9, no. 11, pp. 8515–8530, 2021.
\bibitem{ref17} H. Wang, D. Zhang, Y. Wang, J. Ma, Y. Wang, and S. Li, “Rt-fall: A real-time and contactless fall detection system with commodity wifi devices,” \emph{IEEE Transactions on Mobile Computing}, vol. 16, no. 2, pp. 511–526, 2016.
\bibitem{ref18} L. Zhang, Z. Wang, and L. Yang, “Commercial wi-fi based fall detection with environment influence mitigation,” in \emph{2019 16th Annual IEEE International Conference on Sensing, Communication, and Networking (SECON)}. IEEE, 2019, pp. 1–9.
\bibitem{ref19} Y. Zheng, Y. Zhang, K. Qian, G. Zhang, Y. Liu, C. Wu, and Z. Yang, “Zero-effort cross-domain gesture recognition with wi-fi,” in \emph{Proceedings of the 17th Annual International Conference on Mobile Systems, Applications, and Services}, 2019, pp. 313–325.
\bibitem{ref20} S. Ding, Z. Chen, T. Zheng, and J. Luo, “Rf-net: A unified meta-learning framework for rf-enabled one-shot human activity recognition,” in \emph{Proceedings of the 18th Conference on Embedded Networked Sensor Systems}, 2020, pp. 517–530.
\bibitem{ref21} A. Romero, N. Ballas, S. E. Kahou, A. Chassang, C. Gatta, and Y. Bengio, “Fitnets: Hints for thin deep nets,” \emph{arXiv preprint arXiv:1412.6550}, 2014.
\bibitem{ref22} Y. Wang, S. Yang, F. Li, Y. Wu, and Y. Wang, “Fallviewer: A fine-grained indoor fall detection system with ubiquitous wi-fi devices,” \emph{IEEE Internet of Things Journal}, 2021.
\bibitem{ref23} E. E. Stone and M. Skubic, “Fall detection in homes of older adults using the microsoft kinect,” \emph{IEEE journal of biomedical and health informatics}, 2014.
\bibitem{ref24} J. Wang, Z. Zhang, B. Li, S. Lee, and R. S. Sherratt, “An enhanced fall detection system for elderly person monitoring using consumer home networks,” \emph{IEEE transactions on consumer electronics}, 2014.
\bibitem{ref25} P. Pierleoni, A. Belli, L. Palma, M. Pellegrini, L. Pernini, and S. Valenti, “A high reliability wearable device for elderly fall detection,” \emph{IEEE Sensors Journal}, 2015.
\bibitem{ref26} Q. T. Huynh, U. D. Nguyen, L. B. Irazabal, N. Ghassemian, and B. Q. Tran, “Optimization of an accelerometer and gyroscope-based fall detection algorithm,” \emph{Journal of Sensors}, vol. 2015, 2015.
\bibitem{ref27} Y. Zhang, W. Hou, Z. Yang, and C. Wu, “Vecare: Statistical acoustic sensing for automotive in-cabin monitoring,” in \emph{20th USENIX Symposium on Networked Systems Design and Implementation (NSDI 23)}, 2023, pp. 1185–1200.
\bibitem{ref28} W. Wang, A. X. Liu, and M. Shahzad, “Gait recognition using wifi signals,” in \emph{Proceedings of the 2016 ACM International Joint Conference on Pervasive and Ubiquitous Computing}, 2016, pp. 363–373.
\bibitem{ref29} X. Li, D. Zhang, Q. Lv, J. Xiong, S. Li, Y. Zhang, and H. Mei, “Indotrack: Device-free indoor human tracking with commodity wi-fi,” \emph{Proceedings of the ACM on Interactive, Mobile, Wearable and Ubiquitous Technologies}, 2017.
\bibitem{ref30} K. Qian, C. Wu, Z. Yang, Y. Liu, and K. Jamieson, “Widar: Decimeter-level passive tracking via velocity monitoring with commodity wi-fi,” in \emph{Proceedings of the 18th ACM International Symposium on Mobile Ad Hoc Networking and Computing}, 2017, pp. 1–10.
\bibitem{ref31} F. Zhang, C. Chen, B. Wang, and K. R. Liu, “Wispeed: A statistical electromagnetic approach for device-free indoor speed estimation,” \emph{IEEE Internet of Things Journal}, vol. 5, no. 3, pp. 2163–2177, 2018.
\bibitem{ref32} C. Wu, F. Zhang, Y. Hu, and K. R. Liu, “Gaitway: Monitoring and recognizing gait speed through the walls,” \emph{IEEE Transactions on Mobile Computing}, vol. 20, no. 6, pp. 2186–2199, 2020.
\bibitem{ref33} W. Jiang, C. Miao, F. Ma, S. Yao, Y. Wang, Y. Yuan, H. Xue, C. Song, X. Ma, D. Koutsonikolas et al., “Towards environment independent device free human activity recognition,” in \emph{Proceedings of the 24th annual international conference on mobile computing and networking}, 2018, pp. 289–304.
\bibitem{ref34} J. Zhang, Z. Tang, M. Li, D. Fang, P. Nurmi, and Z. Wang, “Crosssense: Towards cross-site and large-scale wifi sensing,” in \emph{Proceedings of the 24th Annual International Conference on Mobile Computing and Networking}, 2018, pp. 305–320.
\bibitem{ref35} Z. Yang, Y. Zhang, K. Qian, and C. Wu, “Slnet: A spectrogram learning neural network for deep wireless sensing,” in \emph{20th USENIX Symposium on Networked Systems Design and Implementation (NSDI 23)}, 2023, pp. 1221–1236.
\bibitem{ref36} C. Li, Z. Cao, and Y. Liu, “Deep ai enabled ubiquitous wireless sensing: A survey,” \emph{ACM Computing Surveys (CSUR)}, vol. 54, pp. 1–35, 2021.
\bibitem{ref37} Y. Zhang, Y. Zheng, G. Zhang, K. Qian, C. Qian, and Z. Yang, “Gaitsense: towards ubiquitous gait-based human identification with wifi,” \emph{ACM Transactions on Sensor Networks (TOSN)}, 2021.
\bibitem{ref38} D. A. Hill, “Electromagnetic fields in cavities: Deterministic and statistical theories,” \emph{IEEE Antennas and Propagation Magazine}, 2014.
\bibitem{ref39} D. Tse and P. Viswanath, \emph{Fundamentals of wireless communication}. Cambridge university press, 2005.
\bibitem{ref40} F. Zhang, C. Wu, B. Wang, H.-Q. Lai, Y. Han, and K. R. Liu, “Widetect: Robust motion detection with a statistical electromagnetic model,” \emph{Proceedings of the ACM on Interactive, Mobile, Wearable and Ubiquitous Technologies}, vol. 3, no. 3, pp. 1–24, 2019.
\bibitem{ref41} A. Vaswani, N. Shazeer, N. Parmar, J. Uszkoreit, L. Jones, A. N. Gomez, Ł. Kaiser, and I. Polosukhin, “Attention is all you need,” \emph{Advances in neural information processing systems}, vol. 30, 2017.
\bibitem{ref42} K. He, X. Zhang, S. Ren, and J. Sun, “Deep residual learning for image recognition,” in \emph{Proceedings of the IEEE conference on computer vision and pattern recognition}, 2016, pp. 770–778.
\bibitem{ref43} J. Devlin, M.-W. Chang, K. Lee, and K. Toutanova, “Bert: Pre-training of deep bidirectional transformers for language understanding,” \emph{arXiv preprint arXiv:1810.04805}, 2018.
\bibitem{ref44} D. Tran, L. Bourdev, R. Fergus, L. Torresani, and M. Paluri, “Learning spatiotemporal features with 3d convolutional networks,” in \emph{Proceedings of the IEEE international conference on computer vision}, 2015.
\bibitem{ref45} F. Tung and G. Mori, “Similarity-preserving knowledge distillation,” in \emph{Proceedings of the IEEE/CVF International Conference on Computer Vision}, 2019, pp. 1365–1374.
\bibitem{ref46} T. Chen and C. Guestrin, “Xgboost: A scalable tree boosting system,” in \emph{Proceedings of the 22nd acm sigkdd international conference on knowledge discovery and data mining}, 2016, pp. 785–794.
\bibitem{ref47} L. Breiman, “Random forests,” \emph{Machine learning}, vol. 45, pp. 5–32, 2001.

\end{thebibliography}

\section*{作者简介（Author Biographies）}
\noindent\textbf{Guoxuan Chi（IEEE Member）}：2019 年获北京邮电大学信息与通信工程学院学士学位，2024 年获清华大学软件学院博士学位。现为清华大学研究助理。研究方向：无线感知（wireless sensing）与移动计算（mobile computing）。

\noindent\textbf{Guidong Zhang（IEEE Member）}：2018 年获中国科学技术大学电子工程与信息科学系学士学位，2023 年获清华大学软件学院博士学位。研究方向：无线感知与移动计算。

\noindent\textbf{Xuan Ding（IEEE Member）}：2008 年获清华大学学士学位，2014 年获清华大学博士学位。现为清华大学软件学院与 BNRist 研究助理教授。研究方向：隐私保护计算（privacy-preserving computing）、区块链（blockchain）、RFID 与无线感知。

\noindent\textbf{Qiang Ma（IEEE Member）}：2009 年获清华大学计算机科学与技术系学士学位，2013 年获香港科技大学计算机科学与工程系博士学位。现为清华大学软件学院助理研究员。研究方向：无线传感器网络（wireless sensor networks）、移动计算与隐私（privacy）。

\noindent\textbf{Zheng Yang（IEEE Fellow，通讯作者）}：2006 年获清华大学计算机科学学士学位，2010 年获香港科技大学计算机科学博士学位。现为清华大学副教授。研究方向：物联网（Internet of Things）与移动计算。国家杰出青年基金获得者；曾获国家自然科学奖二等奖。

\noindent\textbf{Zhenguo Du}：2007 年获东南大学吴健雄学院学士学位，2012 年获中国科学技术大学电子工程与信息科学系博士学位。现就职于华为技术有限公司，从事无线通信、无线感知与无线定位研究。

\noindent\textbf{Houfei Xiao}：2009 年获华中科技大学电子信息与工程学院学士学位，2014 年获北京邮电大学信息光子与光通信国家重点实验室博士学位。现就职于华为技术有限公司，从事无线通信、无线感知与无线定位研究。

\noindent\textbf{Zhuang Liu}：2013 年获长沙理工大学自动化学士学位，2016 年获华中科技大学电子工程硕士学位。现就职于华为技术有限公司，从事无线通信、无线感知与无线定位研究。

\bigskip
\noindent\textbf{作者照片占位：}（原文含作者照片，请在此处插入）
\end{document}
